\documentclass[11pt,a4paper]{article}
\usepackage{geometry}
\geometry{left=2.5cm,right=2.5cm,top=2.5cm,bottom=2.5cm}
\usepackage{booktabs}
\usepackage{xcolor}
\usepackage{tabularx}
\usepackage{hyperref}
\usepackage{amssymb}

\hypersetup{
    colorlinks=true,
    linkcolor=blue,
    filecolor=magenta,      
    urlcolor=cyan,
}

\definecolor{primary}{HTML}{1A365D}
\definecolor{secondary}{HTML}{2B6CB0}

\usepackage{xepersian}
\settextfont{Tahoma}
\setlatintextfont{Arial}

\begin{document}

\begin{center}
    \vspace*{1cm}
    {\Huge\bfseries طرح پیشنهادی فنی و مالی فاز دوم توسعه فروشگاه ناپ‌کامرس} \\[0.6cm]
    {\Large\bfseries مستندات مالی، وضعیت فازهای توافق‌شده و ماژول‌های توسعه جدید} \\[1.5cm]
\end{center}

\tableofcontents
\newpage

\section{مقدمه و خلاصه مدیریتی}
پروژه تحویلی فاز نخست شامل یک نسخه بهینه‌سازی‌شده از سیستم فروشگاهی چندفروشندگی \lr{nopCommerce 4.90.3} به همراه چهار پلاگین اختصاصی (سیستم خرید گروهی، صفحه تخفیف‌های شگفت‌انگیز، سیستم اعلان‌های پیشرفته کاربران و ماژول پیشرفته روش‌های ارسال شرطی) با موفقیت مستقر و آزمایش گردید. شایان ذکر است که کلیه قابلیت‌های توسعه‌یافته در فازهای پیشین، علاوه بر پشتیبانی در اپلیکیشن موبایل، به گونه‌ای طراحی شده‌اند که به طور کامل در نسخه وب موبایل (PWA) نیز قابل استفاده و بهینه‌سازی شده باشند.

این مستند جهت شفاف‌سازی وضعیت مالی پروژه، تبیین مواردی که پیش‌تر توافق شده و بودجه آن‌ها تخصیص یافته است، و همچنین ارائه‌ی چارچوب ماژول‌های جدید درخواستی (تیکتینگ اکوسیستم، وب‌سرویس‌ها و پنل‌های هوش مصنوعی اقتصادی) تنظیم گردیده است تا تصمیم‌گیری در خصوص فاز جدید تسهیل گردد.

\section{وضعیت مالی پروژه و فازهای توافق‌شده (مصوب)}
بر اساس گفتگوهای انجام شده، فرآیند مالی پروژه در سه فاز اولیه توافق شده و بودجه‌ی آن‌ها در سایت تأمین و تخصیص یافته است:

\subsection{فازهای توافق‌شده و تأمین اعتبار شده (مجموعاً ۴۲,۰۰۰,۰۰۰ تومان)}
\begin{enumerate}
    \item \textbf{بخش اول (تعهدات اولیه فاز ۱):} طراحی سیستم اعلان پیشرفته، صفحه تخفیفات شگفت‌انگیز و سیستم خرید گروهی به ارزش \textbf{۲۲,۰۰۰,۰۰۰ تومان} (تخصیص‌یافته).
    \item \textbf{بخش دوم (روش‌های ارسال فاز ۲):} طراحی و پیاده‌سازی سیستم پیشرفته روش‌های ارسال شرطی چهارگانه \lr{(Courier, Transport, Freight, Express)} به ارزش \textbf{۵,۰۰۰,۰۰۰ تومان} (تخصیص‌یافته).
    \item \textbf{بخش سوم (امکانات هوشمند و امنیت اولیه فاز ۳):} بر اساس توافق مورخ ۲ تیر ۱۴۰۵، بخش‌های زیر با بودجه توافقی \textbf{۱۵,۰۰۰,۰۰۰ تومان} تصویب شده و بودجه پروژه به ۴۲ میلیون تومان افزایش یافته است. پیاده‌سازی این موارد در این فاز تعهد شده است:
    \begin{itemize}
        \item \textbf{ماژول A1 (اتصال به Pyk):} افزودن دکمه لینک اختصاصی در بخش هدر سایت با قابلیت مدیریت آدرس از پنل ادمین.
        \item \textbf{ماژول B (جستجوی هوشمند):} جستجوی صوتی و تصویری هوش مصنوعی (نسخه ابری/محلی).
        \item \textbf{ماژول C (تشخیص محصولات تکراری):} تشخیص محصولات تکراری با هوش مصنوعی جهت جلوگیری از ثبت تکراری کالاها در پنل انبارداری.
        \item \textbf{ماژول D (چت‌بات ادمین):} چت‌بات آنلاین پشتیبانی هوش مصنوعی.
        \item \textbf{ماژول E (پشتیبان‌گیری و بازیابی فروشندگان):} ایجاد امکان پشتیبان‌گیری از کاتالوگ و تصاویر توسط فروشندگان در پنل اختصاصی خود و بازیابی آن پس از تایید مدیریت وب‌سایت.
        \item \textbf{ماژول F (امنیت ورود):} احراز هویت دو مرحله‌ای پیامکی \lr{(SMS 2FA)} برای کاربران و محدودیت ورود به مدیریت ادمین \lr{(IP/Device Binding)}.
    \end{itemize}
\end{enumerate}

\subsection{توضیح در خصوص موارد مشمول در بودجه مصوب فاز ۳}
به استحضار کارفرمای محترم می‌رساند که \textbf{دکمه هدر سایت (ماژول A1 - اتصال به Pyk)، سیستم پشتیبان‌گیری و بازیابی فروشندگان (ماژول E)، و احراز هویت دو مرحله‌ای (ماژول F)} در \textbf{بخش سوم (فاز ۳، سفارش سوم)} و در سقف بودجه توافقی \textbf{۱۵,۰۰۰,۰۰۰ تومان} مصوب‌شده پیش‌تر تعهد گردیده‌اند. این موارد در فایل ارسالی قبلی صرفاً به منظور شفاف‌سازی محدوده کاری و معرفی کامل ماژول‌ها فهرست شده بودند و \textbf{هیچ‌گونه هزینه مضاعف یا محاسبه مجدد} برای آن‌ها در نظر گرفته نشده است. جدول برآورد هزینه‌های جدید صرفاً شامل ماژول‌هایی است که \textbf{کاملاً خارج از محدوده و بودجه فازهای پیشین} قرار دارند.

\section{برآورد مالی ماژول‌ها و نیازمندی‌های جدید (خارج از بودجه مصوب)}
نیازمندی‌های جدید مطرح شده در تاریخ ۹ تیر ۱۴۰۵، خارج از محدوده فازهای قبلی بوده و هزینه‌های پیشنهادی پیاده‌سازی آن‌ها به شرح زیر برآورد می‌گردد. شایان ذکر است که کلیه ماژول‌های جدید زیر، علاوه بر پشتیبانی در اپلیکیشن موبایل بومی، برای استفاده در \textbf{وب‌اپلیکیشن موبایل (Mobile Web App)} نیز بهینه و سازگار طراحی خواهند شد؛ هرچند برخی قابلیت‌ها ذاتاً مختص اپلیکیشن بومی موبایل هستند:

\subsection{۱. سیستم تیکتینگ و چت پیشرفته اکوسیستم \lr{(Laravel + Flutter + Next.js)}}
این سیستم به عنوان یک میکروسرویس مستقل طراحی شده و در تمام بخش‌های اکوسیستم (وب‌سایت پیک، اپلیکیشن پیک، اپلیکیشن انبار و حساب‌های مشتریان) یکپارچه خواهد شد.
\begin{itemize}
    \item \textbf{همگام‌سازی با CRM:} ایجاد پل ارتباطی دوطرفه با \lr{CRM} فعلی شرکت جهت هماهنگی تیکت‌ها از طریق \lr{API}.
    \item \textbf{بستر بلادرنگ (Real-time):} راه‌اندازی ارتباطات چت مستقیم بر پایه \lr{WebSockets}.
    \item \textbf{هزینه پیاده‌سازی:} \textbf{۳۰,۰۰۰,۰۰۰ تومان} (شامل توسعه میکروسرویس بک‌اند، بستر وب‌سوکت، و یکپارچه‌سازی رابط کاربری در اپلیکیشن‌های فلاتر و لندینگ Next.js).
\end{itemize}

\subsection{۲. ماژول تیکتینگ بومی ناپ‌کامرس \lr{(.NET Storefront Ticketing)}}
پیاده‌سازی همین سیستم تیکتینگ به عنوان یک پلاگین بومی در سمت وب‌سایت اصلی ناپ‌کامرس (مستقل از میکروسرویس لاراول) جهت ثبت تیکت و چت آنلاین در حساب کاربری مشتریان ناپ‌کامرس.
\begin{itemize}
    \item \textbf{امکانات:} پنل اختصاصی مشتریان در ناپ‌کامرس، مدیریت تیکت‌های تفکیک‌شده فروشندگان، و چت بلادرنگ با هسته دات‌نت \lr{(SignalR)}.
    \item \textbf{هزینه پیاده‌سازی:} \textbf{۱۲,۰۰۰,۰۰۰ تومان} (شامل افزونه اختصاصی ناپ‌کامرس بدون نیاز به سرور واسط خارجی).
\end{itemize}

\subsection{۳. ماژول G: زیرساخت وب‌سرویس موبایل \lr{(REST Web API)}}
ایجاد لایه وب‌سرویس (\lr{API}) منطبق با ناپ‌کامرس در صورتی که کارفرما تصمیم به ساخت نسخه موبایل کلاینت اختصاصی (\lr{Module A3/A4}) داشته باشد.
\begin{itemize}
    \item \textbf{امکانات:} نگاشت متدها، پیاده‌سازی استاندارد احراز هویت \lr{(JWT)} و پوشش پلاگین‌های اختصاصی (مانند خرید گروهی و تخفیفات شگفت‌انگیز) برای ارتباط با موبایل.
    \item \textbf{هزینه پیاده‌سازی:} \textbf{۱۰,۰۰۰,۰۰۰ تومان}.
\end{itemize}

\subsection{۴. برآورد پنل‌های تجاری هوش مصنوعی اقتصادی}
استفاده از پنل‌های اقتصادی زیر که گزینه‌های بسیار مناسبی برای بازار ایران هستند پیشنهاد می‌گردد:

\paragraph{پنل‌های داخلی پیشنهادی:}
\begin{itemize}
    \item \textbf{پلتفرم \lr{AvalAI} \lr{(avalai.ir)}:} این سرویس دسترسی بدون تحریم و با پرداخت ریالی (شتاب) به مدل‌های متنوع هوش مصنوعی را فراهم می‌کند و پنل مدیریتی دقیقی برای اعمال سقف مصرف روزانه/ماهانه دارد. هزینه توکن‌های آن مطابق با تعرفه جهانی و بسیار ناچیز است.
    \item \textbf{سرویس آوانگار \lr{(Part AI)}:} برای بخش جستجوی صوتی، می‌توانیم از سرویس بومی «آوانگار» استفاده کنیم که برای لهجه‌های فارسی و گفتار بومی بهینه‌سازی شده و هزینه‌ای کاملاً اقتصادی و ریالی دارد.
\end{itemize}

\paragraph{برآورد جدید هزینه‌ها (با فرض مدل‌های اقتصادی):}
\begin{itemize}
    \item \textbf{بخش متنی و چت‌بات پشتیبانی:} با استفاده از مدل‌های سریع و ارزان، هزینه پردازش متن برای ترافیک نرمال سالانه به کمتر از ۵۰۰,۰۰۰ تومان کاهش می‌یابد.
    \item \textbf{بخش چندرسانه‌ای (جستجوی صوتی و تصویری انبار):} با بهینه‌سازی حجم تصاویر ارسالی و محدود کردن زمان صوت، هزینه سالانه این بخش به حدود ۱,۰۰۰,۰۰۰ تا ۳,۰۰۰,۰۰۰ تومان کاهش خواهد یافت.
    \item \textbf{جمع کل برآورد سالانه جدید:} بین ۱.۵ تا ۴ میلیون تومان (بسته به میزان استفاده واقعی کاربران) خواهد بود.
\end{itemize}

\paragraph{نحوه محاسبه هزینه‌ها:}
هزینه استفاده از سرویس‌های هوش مصنوعی به صورت \lr{Token-Based} (پرداخت بر اساس میزان مصرف) محاسبه می‌شود؛ به این معنا که هیچ هزینه ثابت یا اشتراک اجباری برای استفاده از مدل‌ها وجود ندارد و تنها به میزان استفاده واقعی از سرویس‌ها (تعداد درخواست‌ها، حجم متن پردازش‌شده، تصاویر یا فایل‌های صوتی) هزینه محاسبه می‌شود. بنابراین هرچه میزان استفاده کمتر باشد، هزینه نیز به همان نسبت کاهش می‌یابد و در صورت افزایش استفاده، هزینه صرفاً متناسب با میزان مصرف واقعی افزایش خواهد یافت. همچنین با توجه به امکان تعیین سقف بودجه و اعمال محدودیت‌های مصرف در پنل مدیریت، کنترل کامل بر هزینه‌ها وجود خواهد داشت و از هرگونه افزایش ناخواسته هزینه جلوگیری می‌شود.

\subsection{جدول خلاصه برآورد ماژول‌های جدید}
\begin{table}[htbp]
\centering
\begin{tabularx}{\textwidth}{c X c r}
\toprule
\textbf{ردیف} & \textbf{عنوان ماژول جدید} & \textbf{مدت زمان توسعه} & \textbf{هزینه پیشنهادی (تومان)} \\
\midrule
۱ & سیستم تیکتینگ و چت پیشرفته اکوسیستم & ۲ الی ۳ هفته & ۳۰,۰۰۰,۰۰۰ \\
۲ & ماژول تیکتینگ بومی ناپ‌کامرس \lr{(.NET)} & ۱ الی ۲ هفته & ۱۲,۰۰۰,۰۰۰ \\
۳ & ماژول G: زیرساخت وب‌سرویس موبایل & ۱ الی ۲ هفته & ۱۰,۰۰۰,۰۰۰ \\
\bottomrule
\end{tabularx}
\caption{برآورد هزینه‌ها و زمان تحویل ماژول‌های الحاقی جدید فاز دوم --- \textbf{توجه:} ماژول‌های A1، E و F که پیش‌تر در بودجه ۱۵ میلیون تومانی فاز ۳ مصوب شده‌اند در این جدول لحاظ نشده‌اند.}
\end{table}

\subsection{جزئیات و دلایل فنی برآورد هزینه‌ها (توجیه فنی پروژه)}
به منظور شفاف‌سازی و درک بهتر کارفرمای محترم از ابعاد فنی کار، دلایل و چالش‌های پیاده‌سازی هر یک از بخش‌ها به شرح زیر تبیین می‌گردد:

\begin{enumerate}
    \item \textbf{سیستم تیکتینگ و چت پیشرفته اکوسیستم (۳۰,۰۰۰,۰۰۰ تومان):}
    این هزینه به دلیل توسعه یک \textbf{میکروسرویس چندسکویی \lr{(Cross-Platform)} مستقل} است. چالش اصلی این بخش، برنامه‌نویسی و اتصال همزمان ۴ اپلیکیشن مختلف (وب‌سایت پیک با \lr{Next.js}، اپلیکیشن پیک با \lr{Flutter}، اپلیکیشن انبار و پنل مشتریان) به یک پایگاه داده مشترک است. راه‌اندازی سرور اختصاصی وب‌سوکت جهت چت همزمان بلادرنگ بدون تداخل، در کنار پیاده‌سازی ساختار همگام‌سازی اطلاعات تیکت‌ها با سیستم \lr{CRM} فعلی شرکت از طریق وب‌سرویس‌های امن، پیچیدگی فنی بسیار بالایی دارد و نیازمند هماهنگی میان چندین تکنولوژی متفاوت (Laravel, Flutter, Next.js) می‌باشد.
    
    \item \textbf{ماژول تیکتینگ بومی ناپ‌کامرس (۱۲,۰۰۰,۰۰۰ تومان):}
    این ماژول برخلاف ابزارک‌های آماده چت، به صورت یک \textbf{پلاگین بومی دات‌نت \lr{(.NET Core)}} مستقیماً درون ساختار ناپ‌کامرس کدنویسی می‌شود. از مهم‌ترین چالش‌های فنی این بخش، تفکیک دسترسی فروشندگان در سیستم چندفروشندگی است؛ به طوری که هر فروشنده در پنل اختصاصی خود صرفاً تیکت‌های مربوط به محصولات و سفارش‌های خود را ببیند و حریم خصوصی مشتریان حفظ شود. همچنین پیاده‌سازی زیرساخت چت با هسته بومی ناپ‌کامرس با استفاده از فناوری \lr{SignalR} بدون اثر منفی روی کارایی و سرعت لود کل سایت، نیازمند بهینه‌سازی دقیق منابع سرور است.
    
    \item \textbf{ماژول E: پشتیبان‌گیری و بازیابی فروشندگان (مشمول در بودجه مصوب فاز ۳ - بدون هزینه جدید):}
    پشتیبان‌گیری کلی از دیتابیس امری ساده است، اما جداسازی داده‌های مربوط به یک فروشنده خاص در یک سیستم چندفروشندگی بدون افشای داده‌های سایر فروشندگان، چالشی جدی است. سیستم باید فایل‌های دانلودی، ویژگی‌ها، انبار و تصاویر مربوط به همان فروشنده را استخراج کرده و به صورت یک فایل ساختاریافته فشرده \lr{(ZIP)} تحویل دهد. فرآیند بازیابی داده‌ها نیز نیازمند اعتبارسنجی تداخل کدها \lr{(SKU)}، ادغام محصولات ادمین، و سیستم برگشت‌پذیری تراکنش‌ها \lr{(Transactional Rollback)} در صورت بروز هرگونه خطا است تا از تخریب احتمالی دیتابیس کل سایت جلوگیری شود. بر اساس موافقت کارفرما، این ماژول ذیل بخش سوم (امکانات هوشمند و امنیت فاز ۳) و در سقف همان بودجه مصوب ۱۵ میلیون تومانی کل فاز ۳ اجرا شده و شامل هیچ هزینه مازاد جدیدی نمی‌باشد.
    
    \item \textbf{ماژول G: زیرساخت وب‌سرویس موبایل (۱۰,۰۰۰,۰۰۰ تومان):}
    سیستم ناپ‌کامرس به صورت پیش‌فرض فاقد لایه وب‌سرویس امن جهت ارتباط مستقیم با موبایل است. برای این منظور، نیاز به پیاده‌سازی لایه امن وب‌سرویس \lr{(REST API)} با استانداردهای احراز هویت \lr{JWT} و کنترل هدرهای \lr{CORS} است. علاوه بر متدهای استاندارد خرید، چالش اصلی این بخش، نگاشت متدها و قابلیت‌های پلاگین‌های اختصاصی فاز اول (نظیر ماژول خرید گروهی و تخفیفات شگفت‌انگیز) به کدهای وب‌سرویس است تا کلاینت‌های موبایل بتوانند بدون خطا با این امکانات ارتباط برقرار کنند.
\end{enumerate}

\newpage
\section{جزئیات فنی و تصمیمات اتخاذ شده بر اساس بازخورد کارفرما}
جهت تعیین دقیق‌تر معماری و عملکرد ماژول‌ها، بر اساس پاسخ‌های ارائه شده توسط کارفرما، نیازمندی‌ها و جزئیات فنی به شرح زیر نهایی گردید:

\subsection{مشخصات فنی نهایی ماژول E (پشتیبان‌گیری و بازیابی فروشندگان)}
\begin{itemize}
    \item \textbf{ساختار و محتوای پشتیبان:} فایل‌های پشتیبانی که فروشندگان از پنل اختصاصی خود تهیه می‌کنند، دقیقاً باید شامل همان نوع اطلاعاتی باشد که در فایل‌های پشتیبان پنل مدیریت اصلی وب‌سایت ذخیره می‌شود (مشخصات کالا، ویژگی‌ها، تصاویر و...) با این تفاوت که خروجی فیلتر شده و منحصراً حاوی اطلاعات پنل و محصولات همان فروشنده خواهد بود.
    \item \textbf{ذخیره‌سازی و دسترسی:} فایل‌های پشتیبان روی دیسک محلی سرور ذخیره می‌شوند. همچنین فروشندگان قادر خواهند بود یک نسخه کپی از فایل پشتیبان خود را دانلود و نزد خود نگهداری نمایند.
    \item \textbf{فرآیند بازیابی اطلاعات (Restore):} عملیات بازیابی با خواندن فایل پشتیبان آپلود شده توسط فروشنده آغاز می‌شود. قبل از اعمال نهایی فرآیند بازیابی بر روی دیتابیس، مدیر اصلی وب‌سایت به فایل مربوطه دسترسی داشته و امکان تایید یا رد درخواست بازیابی را دارا خواهد بود.
\end{itemize}

\subsection{مشخصات فنی نهایی ماژول C (تشخیص محصولات تکراری با هوش مصنوعی)}
\begin{itemize}
    \item \textbf{دامنه بررسی تکرار:} معیار جلوگیری از ثبت مجدد محصول تکراری، به صورت سراسری و در سطح کل محصولات کل فروشگاه (Global) خواهد بود.
    \item \textbf{عوامل حساسیت هوش مصنوعی:} هسته پردازش تصویر و متن هوش مصنوعی برای تشخیص شباهت کالاها نسبت به موارد زیر حساس خواهد بود:
    \begin{itemize}
        \item نام محصول
        \item تصویر محصول (با پردازش شباهت تصویر)
        \item ویژگی‌ها و مشخصات فنی محصول
        \item توضیحات محصول
    \end{itemize}
    \item \textbf{فرآیند مواجهه با محصول تکراری:} در صورت شناسایی محصول تکراری، از ثبت محصول جدید ممانعت به عمل می‌آید. در عوض، سیستم محصول ثبت‌شده قبلی روی سایت را به فروشنده نمایش می‌دهد تا فروشنده صرفاً بتواند اطلاعات اختصاصی مربوط به فروشگاه خود (نظیر قیمت، موجودی انبار کالا و سایر فیلدهای مجاز که اجازه تغییر آن‌ها را دارد) را برای آن محصول ویرایش و ذخیره کند.
\end{itemize}

\newpage
\section{پیش‌نیازها و سخت‌افزار مورد نیاز سرور هوش مصنوعی محلی}
در صورتی که کارفرما تمایل به پردازش داده‌های هوش مصنوعی روی سرور محلی خود (به‌جای استفاده از پردازش ابری خارجی) داشته باشد، شرایط سخت‌افزاری به شرح زیر پیشنهاد و تحلیل می‌گردد:

\subsection{مشخصات سخت‌افزاری استاندارد سرور ویندوزی}
\begin{itemize}
    \item \textbf{سیستم‌عامل:} \lr{Windows Server 2022}
    \item \textbf{پردازنده \lr{(CPU)}:} حداقل ۱۲ هسته حقیقی (\lr{Intel Xeon Gold} یا \lr{AMD EPYC})
    \item \textbf{حافظه رم \lr{(RAM)}:} حداقل ۶۴ گیگابایت
    \item \textbf{فضای ذخیره‌سازی:} ۱ ترابایت \lr{NVMe SSD}
    \item \textbf{کارت گرافیک \lr{(GPU)}:} حداقل \lr{NVIDIA RTX 4070 Ti} \lr{(12GB)} یا \lr{RTX 4080} \lr{(16GB)} جهت پردازش سریع مدل‌های صوتی و تصویری محلی.
\end{itemize}

\subsection{تحلیل سخت‌افزارهای جایگزین درخواستی کارفرما}
با توجه به هزینه بالای کارت‌های گرافیکی رده مصرف‌کننده، پاسخ به موارد سخت‌افزاری پیشنهادی به شرح زیر است:
\begin{itemize}
    \item \textbf{استفاده از کارت گرافیک \lr{NVIDIA T4}:} 
    استفاده از کارت گرافیک \lr{NVIDIA T4} (دارای ۱۶ گیگابایت حافظه \lr{VRAM}) کاملاً امکان‌پذیر و مناسب است. این کارت برای محاسبات استنتاج هوش مصنوعی در سرورها بهینه‌سازی شده و به دلیل حافظه ۱۶ گیگابایتی، پایداری بسیار خوبی در اجرای مدل‌های صوتی، تصویری و متنی سبک دارد. سرعت پردازش آن نسبت به نسل‌های جدید مانند \lr{RTX 4080} کمتر است اما برای این اکوسیستم کاملاً پاسخگو است.
    \item \textbf{استفاده از ۳۲ گیگابایت حافظه رم \lr{(RAM)}:}
    رم ۳۲ گیگابایت به عنوان حداقل مطلق رم سیستم برای راه‌اندازی مدل‌ها پذیرفتنی است، اما باید در نظر داشت که در صورت اجرای همزمان مدل‌های پردازش زبان طبیعی (چت‌بات)، شباهت تصویر و پردازش گفتار، سیستم ممکن است با کمبود حافظه مواجه شده و به حافظه مجازی \lr{(Swap)} روی دیسک تکیه کند که باعث افت عملکرد می‌شود.
    \item \textbf{استفاده از هارد دیسک معمولی \lr{(HDD)}:}
    استفاده از هارد دیسک معمولی (\lr{HDD}) برای میزبانی مدل‌های هوش مصنوعی به هیچ وجه توصیه نمی‌شود. مدل‌های هوش مصنوعی در قالب فایل‌های چند گیگابایتی ذخیره می‌شوند که برای اجرا باید در حافظه لود شوند. لود اولیه مدل از روی \lr{HDD} به حافظه گرافیک می‌تواند چند دقیقه زمان ببرد. همچنین دیتابیس‌های برداری برای جستجوی صوتی و تصویری تکراری‌ها به سرعت بالایی در خواندن اطلاعات نیاز دارند. توصیه اکید می‌شود حداقل یک درایو \lr{SSD} (حداقل با ظرفیت ۲۵۶ یا ۵۱۲ گیگابایت) برای سیستم‌عامل و قرارگیری فایل مدل‌های هوش مصنوعی تعبیه شده و از \lr{HDD} صرفاً برای ذخیره تصاویر عادی و فایل‌های پشتیبان استفاده گردد.
\end{itemize}

\end{document}
